\section{Introduction}
\label{sec:intro}

%TL;DR (X) Generating buildings in block worlds is a standing 3D generation task, and trained generators need a corpus and a training run for every new condition.
Block worlds such as Minecraft hold millions of player-made buildings, and generating new ones is
a standing task for procedural content and 3D generation~\citep{salge2018gdmc,chen2019voxelcnn,merino2026dreamcubed}.
A building in such a world is a voxel grid in which every cell holds one of hundreds of block
types. Trained generators write this grid cell by cell or denoise it as a
volume~\citep{chen2019voxelcnn,merino2026dreamcubed,pun2025brickgpt}, and each new form of
conditioning, such as a text prompt or a reference image, requires a labelled corpus and a new
training run.

%TL;DR (X) A language model that writes a short build program instead gets text conditioning for free, at a few cents per house.
A language model that writes a build program avoids both. The program is a few dozen commands
in the style of WorldEdit~\citep{worldedit}, such as \texttt{fill}, \texttt{walls} and
\texttt{gable}, and one command places tens to hundreds of blocks. The prompt is the conditioning,
and the model needs no Minecraft training data. With \gptmini{}~\citep{openai2025gpt5} as the
builder, one house on a $32^3$ canvas costs 4.0 to 4.8 US cents in our runs (\cref{sec:method}).
% 0.040 / 0.045 / 0.048 USD per build: card run_20261109_121625 cost table; T51
Several systems already let language models build in Minecraft from text or from a given
reference image~\citep{hu2024t2bm,chen2024apt,fu2026craftutopia,wei2025mineanybuild}.

%TL;DR (X) A popular upgrade for program-writing agents is visual imagination: draw the target image, then build toward it with render-and-compare repair.
A popular upgrade for such agents is visual imagination. A text-to-image model first draws what
the result looks like, and the agent then acts toward that picture. Robot policies
condition on generated subgoal images or videos~\citep{black2024susie,du2023unipi}, web agents
simulate the outcome of an action before taking it~\citep{gu2024webdreamer}, and code-writing
3D agents render their program, compare the render with a target image and edit the
code~\citep{yin2026viga,li2026legoanything}. For a building the appeal is direct: a picture fixes
the massing, the materials and the roof at once, which a short prompt such as
``Small Brick Library (a Minecraft brick house)'' leaves open.

%TL;DR (Y) Measuring what the image adds is hard: an imagination loop also adds repair turns and extra text, and comparing against a one-shot builder credits the image with their effect.
What the image adds is harder to measure than it looks. An imagination loop changes three
things at once. It gives the builder a target picture; it adds repair turns in which the builder
sees its own render; and, in most designs, it adds text about the picture that the builder reads.
A comparison against a builder that writes one program and stops credits the image with the
effect of all three. The image is also an incomplete target for a house: it shows one exterior
view, while what makes a house a house, rooms, doors and stairs a player walks through, sits
behind its walls.

%TL;DR (Z) We build a recipe harness in which every arm differs from its control in one ingredient, with images shared across arms through a content-addressed cache.
We therefore run imagination as a controlled study. A recipe harness defines each arm as one
YAML file over the same builder, build language, renderer and budget, and a content-addressed
image cache gives every arm that draws a dream exactly the same images (\cref{sec:method}).
\classicm{} runs the same number of repair turns and the same best-of selection as the
imagination loop, but critiques its own renders without a target image. \dream{} adds the target
image. \dreamcap{} adds a structured caption of that image, written by a captioner that never sees
the prompt. Three controls separate the caption's content from its length and from the image
itself: \shufcap{} shows the caption of another prompt's dream, \textspec{} shows a specification
written from the prompt with no image, and \bestoffour{} spends the same compute on independent
samples without any repair.

%TL;DR (1a) Finding on appearance: the caption nearly halves the multi-view render distance to real houses, and the image without the caption ties the no-image loop.
On 133 house prompts, the caption carries the visual gain. \dreamcap{} reaches a multi-view
render distance to real houses (\mvkid{}) of 0.217, against 0.403 for \classicm{} and 0.342 for
\dream{}, and beats both in Holm-corrected paired tests ($p=0.014$). \dream{} and \classicm{} share
a significance group. The image alone therefore buys nothing a matched loop without an image
does not, and the gain appears once the image is turned into words the builder reads.
\TODO{n=256 numbers; \shufcap{} and \textspec{} decide whether the gain is the caption's content or its length.}
% MV-KID 0.217 / 0.342 / 0.403, p_holm 0.009 both pairs, dream~classic same group: T51 addendum 2; card run_20261109_121625 (groups d / e / ef)

%TL;DR (1b) Finding on structure: no arm moves geometry, the no-image loop builds the most enterable houses, and all arms repeat themselves far more than real houses do.
Structure does not follow appearance. On the voxel-geometry distance (\geomkid{}) no agentic arm
separates from another (\dream{} 17.9, \classicm{} 18.4, \dreamcap{} 15.7), and a 27M-parameter
voxel transformer trained on houses (\voxelar{}) reaches 9.06. A simulated player enters 68\% of
\classicm{} builds, 50\% of \dream{} builds and 51\% of \dreamcap{} builds, against 93\% of real
houses. Within each arm, 21\% to 32\% of builds have a near-copy in the same set, against none
of the real houses, so all three arms fail the near-twin gate of the evaluation protocol we
adopt~\citep{anonymous2027blindness}. None of them copies a training house.
\TODO{n=256: replace \geomkid{}, \enterable{} and \twinrate{} with the n=256 card; check whether the
distinct-prompt card lowers the twin rate.}
% geom_kid 17.882 / 18.351 / 15.712 (groups ef / ef / de), recanon20m_s0 9.063: card run_20261109_121625; 27M params: T118
% enterable 0.68 / 0.50 / 0.51, real 0.93: card run_20261109_121625 (T63 gives 0.95 at n=128)
% twin rate 0.211 / 0.301 / 0.316: card run_20261109_121625; real 0.000 on this card (0.008 at n=256, T91); voxel_dup_rate 0 for all agentic arms: same card

%TL;DR (C) Contributions: a one-ingredient-at-a-time study of imagination, the finding that captions carry the visual gain, and the finding that neither image nor caption reaches structure.
Our contributions are:
\begin{itemize}[leftmargin=*,itemsep=1pt,topsep=2pt]
  \item A controlled study of visual imagination for a program-writing building agent, in which
    each arm differs from its control in one ingredient: turns, image, caption, caption
    content, or sampling (\cref{sec:method}).
  \item Evidence that the visual gain comes from a caption of the imagined image and not from the
    image itself: the caption nearly halves the multi-view render distance to real houses, and the image
    alone ties a compute-matched loop without an image (\cref{sec:results}).
  \item Evidence that neither the image nor the caption reaches structure: geometry is not
    separated, the loop without an image builds the most enterable houses, and every arm repeats
    near-identical builds far more often than real houses do (\cref{sec:results}).
\end{itemize}
